\chapter{Requirement Analysis}

The objective of the requirement analysis is to establish clear expectations of the PCMA's joint design for the purposes of concept development.

\section{Identification of Stakeholders}

The first step in the requirement analysis process is to identify stakeholders from the operational and support life cycle phases of the project. Stakeholders will be classified into the following categories, as defined below:
\begin{itemize}
	\item Primary stakeholders directly interact with the system while it is in operation.
	\item Secondary stakeholders contribute to the development and functioning of the system.
	\item Indirect stakeholders rely on or are affected by the system, but do not interact with the system directly.
\end{itemize}
The stakeholders relevant to the project are defined in Table~\ref{tab:SID table} below.
\begin{table}[htbp]
	\centering
	\caption{Identified stakeholders and their classifications}
	\label{tab:SID table}
	\begin{tabular}{|l|l|l|}
		\hline
		\textbf{Stakeholder ID} & \textbf{Description} & \textbf{Classification} \\ \hline
		SID\_01 & Supervisor of the project & Primary \\ \hline
		SID\_02 & Students researching PCMA development & Secondary \\ \hline
		SID\_03 & Companies invested in PCMA development & Indirect\\ \hline
		SID\_04 & Mechanical workshop staff & Indirect \\ \hline
		SID\_05 & Mechatronics laboratory managers & Indirect \\ \hline
	\end{tabular}
\end{table}

\section{System Context}

The next step is 